\DocumentMetadata{
  pdfversion=2.0,
  pdfstandard=ua-2,
  lang=en-US,
  tagging=on,
  tagging-setup={math/setup=mathml-SE,tabsorder=structure}
}
\documentclass[11pt,letterpaper]{article}
\usepackage[margin=0.68in,includefoot]{geometry}
\usepackage{newcomputermodern}
\usepackage{amsmath,amscd,mathtools}
\providecommand{\square}{\mdlgwhtsquare}
\usepackage[hidelinks]{hyperref}
\hypersetup{
  pdftitle={Numerical Analysis First-Year Exam, May 2022},
  pdfauthor={University of Florida Department of Mathematics},
  pdfsubject={Graduate mathematics examination},
  pdfdisplaydoctitle=true,
  bookmarksopen=true
}
\setcounter{secnumdepth}{0}
\setlength{\parindent}{0pt}
\setlength{\parskip}{0.28em}
\setlength{\emergencystretch}{3em}
\setlength{\leftmargini}{2.15em}
\setlength{\leftmarginii}{2.65em}
\setlength{\leftmarginiii}{3em}
\renewcommand{\labelenumi}{\textbf{\theenumi.}}
\pagestyle{empty}
\raggedbottom
\allowdisplaybreaks
\newenvironment{examproblems}[1]{%
  \begin{enumerate}%
  \setcounter{enumi}{#1}%
  \setlength{\topsep}{0.35em}%
  \setlength{\partopsep}{0pt}%
  \setlength{\itemsep}{0.48em}%
  \setlength{\parsep}{0.18em}%
}{\end{enumerate}}
\newenvironment{examsubparts}{%
  \begin{enumerate}%
  \setlength{\topsep}{0.18em}%
  \setlength{\partopsep}{0pt}%
  \setlength{\itemsep}{0.16em}%
  \setlength{\parsep}{0.1em}%
}{\end{enumerate}}
\newenvironment{examinstructions}{%
  \par\begingroup\itshape
}{%
  \par\endgroup\vspace{0.18em}
}
\makeatletter
\renewcommand\section{\@startsection{section}{1}{0pt}%
  {-1.25ex plus -.25ex minus -.1ex}%
  {0.55ex plus .1ex}%
  {\normalfont\large\bfseries}}
\renewcommand{\@maketitle}{%
  \newpage\null\vskip -1.2em%
  {\footnotesize\bfseries
    \href{https://www.ufl.edu/}{University of Florida}, \href{https://math.ufl.edu/}{Department of Mathematics}\par}%
  \vskip 0.18em%
  \tagstructbegin{tag=Title}%
    {\LARGE\bfseries\href{https://gma.math.ufl.edu/exams/numerical-analysis-first-year/}{Numerical Analysis First-Year Exam}, May 2022\par}%
  \tagstructend%
  \vskip 0.35em%
  \tagmcbegin{artifact}%
    \hrule height 1pt%
  \tagmcend%
  \vskip 0.55em%
}
\makeatother
\title{Numerical Analysis First-Year Exam, May 2022}
\author{}
\date{}
\begin{document}
\maketitle
\thispagestyle{empty}
\begin{examinstructions}
Do 4 (four) problems.
\end{examinstructions}
\begin{examproblems}{0}
\item
This problem has two parts.
\begin{examsubparts}
\item[\textnormal{(a)}]
Let \(G=[0,2]\) and 
\[
g(x)=\frac{1}{3}\left(\frac{x^3}{3}-x^2-\frac{5}{4}x+4\right).
\]
 Use the contraction mapping theorem to prove that if \(x_0\in G\), then the sequence defined by \(x_{k+1}=g(x_k)\), \((k=0,1,\ldots)\) converges to a unique fixed point \(z\in G\).
\item[\textnormal{(b)}]
Consider the fixed point iteration method \(x_{k+1}=g(x_k)\), \(k=0,1,\ldots\) for solving the nonlinear equation \(f(x)=0\). Consider choosing an iteration function of the form 
\[
g(x)=x-af(x)-b(f(x))^2-c(f(x))^3,
\]
 where~\(a\), \(b\), and~\(c\) are parameters to be determined. Assume~\(f\) is sufficiently differentiable and the corresponding iterations \(x_{k+1}=g(x_k)\) converge to a unique fixed point~\(z\). Find expressions for the parameters~\(a\), \(b\), and~\(c\) in terms of functions of~\(z\) such that the iteration method is of fourth order.
\end{examsubparts}
\item
Consider \(f(t)=\sin(t)\).
\begin{examsubparts}
\item[\textnormal{(a)}]
Without using orthogonal polynomials, find the best approximation \(p_1(t)\in \mathbb{P}^2[-1,1]\) to \(f(t)\in C[-1,1]\) with respect to the \(L^2\) norm.
\item[\textnormal{(b)}]
Find the Taylor polynomial approximation \(p_2(t)\) of degree 3 at \(t=0\).
\item[\textnormal{(c)}]
Find the Lagrange polynomial approximation \(p_3(t)\) of degree 3 that interpolates \(f(t)\) at \(t=-1,-\frac{1}{3},\frac{1}{3},1\).
\end{examsubparts}
\item
Given a differentiable function \(f(x)\), consider the problem of finding a polynomial \(p(x)\in\mathbb{P}^n\) such that 
\[
p(x_0)=f(x_0),\qquad p'(x_i)=f'(x_i),\qquad i=1,2,\ldots,n,
\]
 where \(x_i\), \(i=1,2,\ldots,n\), are distinct nodes. (It is not excluded that \(x_1=x_0\).) Show that the problem has a unique solution and describe how it can be obtained.
\item
Let \(Q_m(f)\) denote the~\(m\)-point Gaussian quadrature rule over the interval \([a,b]\) and with continuous weight function \(\rho(x)\geq 0\), that is, 
\[
Q_m(f)=\sum_{i=1}^m\alpha_i f(x_i)\approx J(f)=\int_a^b\rho(x)f(x)\,dx.
\]
 Show that, if~\(a\) and~\(b\) are finite and~\(f\) is continuous, then \(Q_m(f)\to J(f)\) as \(m\to\infty\).
\item
Consider numerically solving the initial value problem 
\[
y'(t)=f(t,y),\quad 0<t\leq t_f,\qquad \text{with }y(0)=\eta.
\]
 Assume~\(f\) is sufficiently differentiable and let~\(h\) denote the step size. Show that all convergent members of the family of methods 
\[
y_{n+2}+(\theta-2)y_{n+1}+(1-\theta)y_n=\frac{1}{4}h\big[(6+\theta)f_{n+2}+3(\theta-2)f_n\big],
\]
 parameterized by~\(\theta\), are also \(A_0\)-stable.
\end{examproblems}
\end{document}
